% ==========================================
% CHAPTER 03: Levels 4 and 5 (Intermediate NodeMCU)
% ==========================================
\chapter{Level 4 \& 5: Rudimentary Obfuscation (NodeMCU)}

\textbf{In This Chapter:}
\begin{itemize}
    \item Level 4: Initializing the native \texttt{server.authenticate} module.
    \item Level 5: Exploring basic bitwise XOR logic for telemetry obfuscation.
\end{itemize}

Prior to analyzing how hackers natively slice HTTP request queries, we will examine two intermediate iterations of the NodeMCU framework (Levels 4 and 5) that attempted to inject primary security via rudimentary code modifications. 

\section{Level 4: The First Password}
In Level 4, the developer decides to stop the public nature of the device by utilizing the built-in \texttt{authenticate} function of the C++ header.

\subsection{NodeMCU Implementation (Level 4)}
This is the unabridged iteration of Level 4. Notice the boolean structure around line 49.

\begin{lstlisting}[language=C++, caption=23022026\_Level\_4\_Password.ino]
/*
  NodeMCU (ESP8266) + DHT22 Secure WiFi Monitor
  ---------------------------------------------
  What this program does:
  1. Connects NodeMCU to your WiFi network
  2. Reads Temperature and Humidity from DHT22 sensor (connected to D5)
  3. Creates a password-protected webpage
  4. Shows values in Cyan (Temperature) and Pink (Humidity)
  5. Automatically refreshes every 5 seconds
  6. Only users with correct username and password can view data
*/

#include <ESP8266WiFi.h>        // Library for WiFi connection
#include <ESP8266WebServer.h>   // Library to create web server
#include <DHT.h>                // Library to use DHT sensor

#define D5PIN D5               // Define D5 pin for DHT22 data

// ----------- WiFi Credentials -----------
const char* ssid = "YOUR_WIFI";
const char* pass = "YOUR_WIFI_PASSWORD";

// ----------- Web Login Credentials -----------
const char* www_username = "admin";   // Change username if needed
const char* www_password = "1234";    // Change password if needed

DHT dht(D5PIN, DHT22);          // Create DHT object (Pin, Sensor type)
ESP8266WebServer server(80);    // Create web server on port 80

void setup() {
  Serial.begin(9600);           // Start serial communication (to show IP)
  dht.begin();                  // Start DHT sensor

  WiFi.begin(ssid, pass);       // Start connecting to WiFi
  while (WiFi.status() != WL_CONNECTED) {
    delay(500);                 // Wait until WiFi connects
  }

  Serial.println("Connected to WiFi");
  Serial.print("Open this IP in browser: ");
  Serial.println(WiFi.localIP());   // Print IP address

  // When someone opens the IP address
  server.on("/", []() {

    // ---- Authentication Check ----
    // If username/password incorrect, ask for login
    if (!server.authenticate(www_username, www_password)) {
      return server.requestAuthentication();
    }

    float t = dht.readTemperature();   // Read temperature
    float h = dht.readHumidity();      // Read humidity

    // ---- Create Simple HTML Page ----
    String page = "<html>";
    page += "<head>";
    page += "<meta http-equiv='refresh' content='5'>";  // Auto refresh every 5 sec
    page += "</head>";
    page += "<body bgcolor='black' text='white'>";
    page += "<center>";

    page += "<h1>Temperature</h1>";
    page += "<h1><font color='cyan'>" + String(t) + " C</font></h1>";

    page += "<h1>Humidity</h1>";
    page += "<h1><font color='pink'>" + String(h) + " %</font></h1>";

    page += "</center>";
    page += "</body>";
    page += "</html>";

    server.send(200, "text/html", page);  // Send webpage to browser
  });

  server.begin();   // Start the web server
}

void loop() {
  server.handleClient();   // Continuously handle browser requests
}
\end{lstlisting}

\section{Level 5: Encryption-Only Monitor}
Level 5 is an educational iteration demonstrating how physical microcontrollers process mathematics. The developer uses an XOR (Exclusive OR) logical bitwise operator (`^`) to mathematically mutate the temperature metric before it flies over the Wi-Fi. 

The HTML JavaScript then executes the same reverse XOR operation natively inside the Chrome engine using the shared integer `123`.

\subsection{NodeMCU Implementation (Level 5)}
\begin{lstlisting}[language=C++, caption=23022026\_Level\_5\_Encryption.ino]
/*
  NodeMCU (ESP8266) + DHT22
  Encryption-Only WiFi Monitor
  --------------------------------
  What this program does:
  1. Connects NodeMCU to WiFi
  2. Reads Temperature & Humidity from DHT22 (D5)
  3. Encrypts values using XOR with a secret key
  4. Sends encrypted values to browser
  5. Browser decrypts and displays real values
  6. Auto refresh every 5 seconds

  NOTE:
  This is educational encryption only (not HTTPS).
*/

#include <ESP8266WiFi.h>
#include <ESP8266WebServer.h>
#include <DHT.h>

#define D5PIN D5

// Replace with your WiFi credentials
const char* ssid = "YOUR_WIFI";
const char* pass = "YOUR_PASSWORD";

DHT dht(D5PIN, DHT22);
ESP8266WebServer server(80);

int secretKey = 123;   // Encryption key (change if needed)

void setup() {
  Serial.begin(9600);
  dht.begin();

  WiFi.begin(ssid, pass);
  while (WiFi.status() != WL_CONNECTED) {
    delay(500);
  }

  Serial.print("Open this IP in browser: ");
  Serial.println(WiFi.localIP());

  server.on("/", []() {

    float t = dht.readTemperature();   // Read temperature
    float h = dht.readHumidity();      // Read humidity

    // Encrypt values (multiply by 10 to preserve decimal)
    int encT = ((int)(t * 10)) ^ secretKey;
    int encH = ((int)(h * 10)) ^ secretKey;

    String page = "<html><head>";
    page += "<meta http-equiv='refresh' content='5'>";
    page += "</head><body bgcolor='black' text='white'><center>";

    page += "<h1>Encrypted Sensor Monitor</h1>";

    // JavaScript decrypts values in browser
    page += "<script>";
    page += "var key=" + String(secretKey) + ";";
    page += "var t=(" + String(encT) + "^key)/10;";
    page += "var h=(" + String(encH) + "^key)/10;";
    page += "document.write('<h1 style=\"color:cyan\">'+t+' C</h1>');";
    page += "document.write('<h1 style=\"color:pink\">'+h+' %</h1>');";
    page += "</script>";

    page += "</center></body></html>";

    server.send(200, "text/html", page);
  });

  server.begin();
}

void loop() {
  server.handleClient();
}
\end{lstlisting}

\begin{takeawaybox}
While Level 5 demonstrates symmetric logic, transmitting the mathematical decryption key natively inside the HTML payload (\texttt{var key=123}) completely nullifies the security. Modern systems rely on TLS certificates (HTTPS) explicitly computed via extreme mathematics to generate secure tunnels, not mere XOR obfuscation. Let's move to Level 6 to explore how injection bypasses authentication entirely.
\end{takeawaybox}
